% =====================================================================
%  PKI Session 9: Capstone Defense Phase 3 & Course Synthesis
%  Author: Nguyen An Khuong (CSE-HCMUT · ORCID: 0000-0002-9910-6387)
%  Teaching Assistant: Nguyen Huu Hoang Huy (LinkedIn)
%  Master of Science in Information Security (CHAT4P) · HK261
% =====================================================================

\title{\color{hcmutblue}\textbf{\LARGE PUBLIC KEY INFRASTRUCTURE}\\[0.25em]
       \color{hcmutblue}\textbf{\large S09: Báo Cáo Đồ Án (Đợt 3) \& Tổng Kết Học Phần}}
\date{\footnotesize Session 9 of 9 $\cdot$ ATHK 512 (KMA CHAT4P, HK261)\\
      Class: Sat 28/11/2026 $\cdot$ 13h30--18h30\\
      \textit{Synthesis: Nhóm 5 (TYN), Đối soát Bảng điểm KMA, Bế mạc}}

\begin{document}
\PKItitleframe

\section{Tổng quan \& Lộ trình Buổi S09}

\begin{frame}{Lộ Trình Học Phần (Buổi 9/9 -- Chung Cuộc)}
  \centering
  \makebox[\linewidth][c]{\PKIcoursemap{9}}
  \vspace{0.4em}
  \raggedright\small
  Buổi học bế mạc học phần ATHK 512: Báo cáo đồ án Nhóm 5 (Certificate Compliance), hệ thống hóa 10 nguyên tắc vàng của PKI và tiến hành đối soát, ký biên bản điểm số chính thức.
  \gvnote{[~2 phút | 13h30--13h32] Chào mừng toàn thể lớp CHAT4P đến với buổi học cuối cùng của học phần.}
\end{frame}

\begin{frame}{Lịch Trình Buổi Bế Mạc S09}
  \begin{columns}[T]
    \begin{column}{0.48\textwidth}
      \textbf{Nửa Đầu Buổi Học (13h30--15h15):}
      \begin{itemize}\small
        \item Báo cáo đồ án Nhóm 5: Certificate Compliance \& CT Logs (Nhóm 2 phản biện).
        \item Đánh giá và tổng kết chất lượng 5 đề tài đồ án chuyên đề.
      \end{itemize}
    \end{column}
    \begin{column}{0.48\textwidth}
      \textbf{Nửa Sau Buổi Học (15h30--18h00):}
      \begin{itemize}\small
        \item Tổng kết 10 nguyên tắc vàng quản trị hạ tầng PKI doanh nghiệp.
        \item Định hướng nghiên cứu luận văn thạc sĩ.
        \item \textbf{Đối soát bảng điểm 100 điểm \& Ký biên bản điểm KMA.}
      \end{itemize}
    \end{column}
  \end{columns}

  \begin{keybox}[Khẳng Định Năng Lực Học Viên]
    Trải qua 9 buổi học và 5 đồ án lớn, học viên lớp CHAT4P đã hoàn thiện trọn vẹn năng lực từ thẩm định toán học, cấu trúc chuẩn nhị phân đến vận hành PKI an ninh mức doanh nghiệp.
  \end{keybox}
  \gvnote{[~3 phút | 13h32--13h35] Chúc mừng học viên đã hoàn thành hành trình 9 buổi học thực chiến.}
\end{frame}

\section{Bảo vệ Đồ án: Nhóm 5 (TYN)}

\begin{frame}{Báo Cáo Đồ Án: Nhóm 5 (Certificate Compliance)}
  \begin{columns}[T]
    \begin{column}{0.5\textwidth}
      \textbf{Đề tài: Kiểm Định Tuân Thủ CT Logs \& ZLint}
      \begin{itemize}\small
        \item \textbf{Thành viên:} Võ Minh Nhật (Trưởng nhóm), Hồ Quốc Thái, Nguyễn Công Ý.
        \item \textbf{Đơn vị phản biện:} Nhóm 2 (RootCA).
        \item \textbf{Mục tiêu kỹ thuật:} Xây dựng hệ thống tự động quét và phân tích chứng chỉ vi phạm CA/B Forum bằng ZLint viết trên nền Golang.
      \end{itemize}
    \end{column}
    \begin{column}{0.46\textwidth}
      \begin{quizbox}[Câu Hỏi Trọng Tâm Phản Biện]
        Làm thế nào hệ thống của Nhóm 5 phân biệt được một lỗi cú pháp ASN.1 vô hại với một lỗ hổng an ninh buộc CA phải thu hồi chứng chỉ trong 24 giờ?
      \end{quizbox}
    \end{column}
  \end{columns}

  \vspace{0.6em}
  \centering\large\sffamily\bfseries\color{hcmutblue}
  Mời đại diện Nhóm 5 lên bục thuyết trình và Live Demo!
  \gvnote{[~45 phút | 13h35--14h20] Phiên báo cáo và phản biện của Nhóm 5.}
\end{frame}

\section{Hệ thống hóa Kiến thức \& 10 Nguyên tắc Vàng}

\begin{frame}{10 Nguyên Tắc Vàng Quản Trị Hạ Tầng PKI Doanh Nghiệp}
  \begin{enumerate}\footnotesize
    \item \textbf{Air-Gapped Root CA:} Root CA tuyệt đối ngắt kết nối mạng, bảo vệ trong HSM FIPS 140-3 L4.
    \item \textbf{Path Length Constraint:} Luôn đặt \texttt{pathLen=0} trên Intermediate CA cấp phát.
    \item \textbf{Khai tử Common Name:} 100\% định danh tên miền bắt buộc phải nằm trong SAN.
    \item \textbf{Chứng chỉ ngắn hạn:} Giảm thời hạn xuống dưới 90 ngày kết hợp ACME RFC 8555 tự động.
    \item \textbf{Khóa chặt Soft-fail:} Kích hoạt OCSP Stapling đi kèm cờ Must-Staple RFC 7633.
    \item \textbf{Minh bạch công khai:} Ghi nhận toàn bộ chứng chỉ vào Certificate Transparency Logs.
    \item \textbf{Zero Trust mTLS:} Xác thực hai chiều cho mọi microservices bằng SPIFFE/SPIRE.
    \item \textbf{Không commit khóa lên Git:} Bắt buộc kiểm soát bằng TruffleHog / Gitleaks pre-commit hooks.
    \item \textbf{Sẵn sàng Hậu lượng tử (PQC-Ready):} Thiết kế kiến trúc lai Hybrid X.509 RFC 9881.
    \item \textbf{Kiểm toán định kỳ:} Chạy ZLint và rà soát CP/CPS RFC 3647 ít nhất một lần mỗi năm.
  \end{enumerate}
  \gvnote{[~5 phút | 14h20--14h25] Đúc kết 10 nguyên tắc quản trị vàng làm hành trang nghề nghiệp cho các tân thạc sĩ An toàn thông tin.}
\end{frame}

\section{Đối soát Bảng điểm \& Bế mạc Học phần}

\begin{frame}{Quy Trình Đối Soát Điểm \& Chốt Điểm KMA}
  \begin{table}[h]
    \centering\scriptsize
    \begin{tabularx}{\linewidth}{@{}p{0.25\linewidth} c p{0.30\linewidth} X@{}}
      \toprule
      \textbf{Thành Phần} & \textbf{Tỷ Lệ} & \textbf{Hình Thức Đánh Giá} & \textbf{Cơ Chế Đối Soát} \\
      \midrule
      \textbf{Chuyên Cần} & 10\% & Điểm danh 9 buổi học & Ký tên đối chiếu tại lớp \\
      \textbf{Thực Hành Lab} & 20\% & 5 bài lab (Log + Git commit) & Đối soát trên repo nhóm \\
      \textbf{Đồ Án Nhóm} & 30\% & Báo cáo Hội đồng + Chấm chéo & Công bố điểm phản biện \\
      \textbf{Thi Cuối Kỳ} & 40\% & Bài thi KTHP tổng hợp & Thẩm định theo barem KMA \\
      \bottomrule
    \end{tabularx}
  \end{table}

  \begin{block}{Công Thức Điểm Tổng Kết}
    $$\text{Điểm TK} = 0.1 \times \text{CC} + 0.2 \times \text{Lab} + 0.3 \times \text{Đồ\_Án} + 0.4 \times \text{Thi}$$
  \end{block}
  \gvnote{[~5 phút | 14h25--14h30] Phổ biến quy chế đối soát điểm công khai, minh bạch trước toàn thể học viên.}
\end{frame}

\begin{frame}{Lời Cảm Ơn \& Bế Mạc Học Phần}
  \centering
  {\Huge\sffamily\bfseries\color{hcmutblue} CHÚC MỪNG KHÓA HỌC CHAT4P!}\\[0.4cm]
  {\Large\sffamily\color{solgreen} Hoàn thành xuất sắc học phần Cơ sở hạ tầng khóa công khai (ATHK 512)}\\[1.2cm]

  \begin{columns}[T]
    \begin{column}{0.48\textwidth}
      \centering
      \textbf{Giảng viên phụ trách:}\\
      Nguyễn An Khương\\
      {\footnotesize\textit{Khoa Khoa học \& Kỹ thuật Máy tính}}\\
      {\footnotesize\textit{Trường ĐHBK Tp. HCM --- ĐHQG-HCM}}\\
      {\scriptsize ORCID: \href{https://orcid.org/0000-0002-9910-6387}{0000-0002-9910-6387}}
    \end{column}
    \begin{column}{0.48\textwidth}
      \centering
      \textbf{Trợ giảng học phần:}\\
      Nguyễn Hữu Hoàng Huy\\
      {\footnotesize\textit{Hỗ trợ kỹ thuật \& Lab Framework}}\\
      {\scriptsize LinkedIn: \href{https://www.linkedin.com/in/nguyen-huu-hoang-huy-abb138366/}{Profile}}
    \end{column}
  \end{columns}

  \vspace{1.0cm}
  {\small\itshape\color{gray} Thành phố Hồ Chí Minh --- Tháng 11 năm 2026}
  \gvnote{[~5 phút | 14h30--14h35] Lời cảm ơn bế mạc và chụp ảnh lưu niệm cùng lớp học.}
\end{frame}

\end{document}
